\section{NHNProvider}
\label{sec:NHN-provider}
\texttt{NHNVideoProvider} oppretter rom som bookinger i NHN sitt API. I \texttt{CreateRoomAsync}, vist i kodeliste~\ref{lst:nhn-create-room}, settes \texttt{startTime} til gjeldende UTC-tid, mens \texttt{endTime} settes til samme tidspunkt pluss \texttt{MeetingDurationHours} fra konfigurasjonen. Romnavnet benyttes som \texttt{subject}, og \texttt{TemplateId} sendes som møtemal. Etter \texttt{POST}-kallet til \texttt{api/client/v1/booking/} tilordnes booking-ID, emne, status og URL-er til \texttt{VideoRoom}.

\begin{lstlisting}[language={[Sharp]C}, label=lst:nhn-create-room, caption={Opprettelse av NHN Booking Request.}]
// HDO-VideoAPI/Providers/NHN/NHNVideoProvider.cs
public async Task<VideoRoom> CreateRoomAsync(CreateRoomRequest request, CancellationToken ct = default)
{
    // ...
    
    var now = DateTime.UtcNow;
    var body = new CreateNHNBookingRequest
    {
        StartTime = now.ToString("yyyy-MM-ddTHH:mm:ssZ"),
        EndTime = now.AddHours(_config.MeetingDurationHours).ToString("yyyy-MM-ddTHH:mm:ssZ"),
        Subject = request.Name ?? Guid.NewGuid().ToString(),
        Template = _config.TemplateId   // Møtemal fra HDO
    };

    // ...

    var response = await httpClient.PostAsJsonAsync("api/client/v1/booking/", body, ct);
    var booking = await response.Content.ReadFromJsonAsync<NHNBooking>(cancellationToken: ct);
    return MapToVideoRoom(booking);
}
\end{lstlisting}

% WebRTC-lenker håndteres enklere enn hos AntMedia. \texttt{GenerateWebRtcTokenAsync} utfører ikke et separat tokenkall, men henter i stedet bookingen og returnerer \texttt{GuestUrl} for avspilling (play) eller \texttt{HostUrl} for øvrige roller. Siden NHN ikke tilbyr tilsvarende mekanisme for tokenutløp, benytter implementasjonen en fast placeholder-verdi. Før URL-ene returneres, transformeres de via \texttt{RewriteNhnUrl}, som omskriver NHN-lenker fra \texttt{/conference?id\&pin=...} til \texttt{?conference=id\&pin=...}. Dette er nødvendig fordi klienten forventer konferanse-ID som en standard query-parameter.

WebRTC-lenker håndteres enklere enn hos AntMedia. \texttt{GenerateWebRtcTokenAsync}, vist i kodeliste~\ref{lst:nhn-webrtc}, utfører ikke et separat tokenkall, men henter bookingen og returnerer \texttt{GuestUrl} for avspilling eller \texttt{HostUrl} for øvrige roller. Siden NHN ikke tilbyr tokenutløp, benyttes en fast placeholder-verdi. Før URL-ene returneres, transformeres de via \texttt{RewriteNhnUrl}, som omskriver NHN-lenker fra \texttt{/conference?id\&pin=...} til \texttt{?conference=id\&pin=...} slik at klienten får konferanse-ID som en standard query-parameter.

\begin{lstlisting}[language={[Sharp]C}, label=lst:nhn-webrtc, caption={Opprettelse av NHN WebRTC-lenke.}]
// HDO-VideoAPI/Providers/NHN/NHNVideoProvider.cs
public async Task<WebRtcLinkResponse> GenerateWebRtcTokenAsync(string roomId, string type, CancellationToken ct = default)
{
    var room = await GetRoomAsync(roomId, ct); // Hent booking fra NHN
    var tokenId = type == "play" 
        ? (room.GuestUrl ?? string.Empty)                   // Gjest-URL for avspilling
        : (room.HostUrl ?? room.GuestUrl ?? string.Empty);  // Vert-URL, fallback til GuestUrl
    return new WebRtcLinkResponse(tokenId, roomId, type, NoExpiry);
}
\end{lstlisting}

% Autentiseringen mot NHN håndteres i \texttt{NHNAuthHandler}. Handleren henter bearer-token fra \texttt{/Token} ved bruk av password grant, basert på brukernavn og passord fra konfigurasjonen. Tokenet lagres i \texttt{NHNTokenCache}, slik at flere handler-instanser kan gjenbruke samme token (jf. kodeliste~\ref{lst:nhn-token-cache}). Cachen tar hensyn til \texttt{expires\_in} fra NHN og trekker fra fem minutter som sikkerhetsmargin for å unngå at et token brukes etter at det har utløpt. I tillegg settes en øvre grense på 13 dager — NHN utsteder tokens med 14 dagers gyldighet, og grensen på 13 dager er differansen etter at sikkerhetsmarginen er trukket fra. Også her benyttes \texttt{SemaphoreSlim} med double-check pattern for å unngå parallelle tokenfornyelser. Feil fra NHN-API-et håndteres av hjelpemetoden \texttt{EnsureSuccessOrThrowAsync}: HTTP-statuskoden \texttt{404 Not Found} kaster \texttt{NotFoundException("Booking", ...)}, mens øvrige feil kaster \texttt{ExternalServiceException("NHN", ...)} med metode, booking-ID, statuskode og responsinnhold.

% \begin{lstlisting}[language={[Sharp]C}, label=lst:nhn-token-cache, caption={NHN token cache logikk}]
% // HDO-VideoAPI/Providers/NHN/NHNAuthHandler.cs
% cache.Token = result!.AccessToken;
% var serverExpiry = DateTime.UtcNow.AddSeconds(
%     result.ExpiresIn - TokenExpiryBufferSeconds);
% cache.Expiry = serverExpiry < DateTime.UtcNow.AddDays(NHNTokenMaxAgeDays)
%     ? serverExpiry
%     : DateTime.UtcNow.AddDays(NHNTokenMaxAgeDays);
% \end{lstlisting}

Autentiseringen mot NHN håndteres i \texttt{NHNAuthHandler}. Handleren henter bearer-token fra \texttt{/Token} ved bruk av password grant, basert på brukernavn og passord fra konfigurasjonen. Tokenet lagres i \texttt{NHNTokenCache} slik at flere handler-instanser kan gjenbruke samme token. Cachen tar hensyn til \texttt{expires\_in} fra NHN og trekker fra fem minutter som sikkerhetsmargin for å unngå at et utløpt token benyttes. I tillegg begrenses levetiden til 13 dager. NHN utsteder tokens med 14 dagers gyldighet, og de 13 dagene er nettoverdien etter at sikkerhetsmarginen er trukket fra. Også her benyttes \texttt{SemaphoreSlim} med double-check-mønsteret (jf.\ kodeliste~\ref{lst:antmedia-auth}) for å unngå parallelle tokenfornyelser. Kodeliste~\ref{lst:nhn-token-cache} viser password grant-requesten og utløpsberegningen, der cache-sjekken er utelatt. Feil fra NHN-API-et håndteres av hjelpemetoden \texttt{EnsureSuccessOrThrowAsync}: HTTP-statuskoden \texttt{404 Not Found} kaster \texttt{NotFoundException("Booking", ...)}, mens øvrige feil kaster \texttt{ExternalServiceException("NHN", ...)} med metode, booking-ID, statuskode og responsinnhold.

\begin{lstlisting}[language={[Sharp]C}, label=lst:nhn-token-cache, caption={Token-henting i NHNAuthHandler}]
// HDO-VideoAPI/Providers/NHN/NHNAuthHandler.cs
private async Task<string> GetTokenAsync(CancellationToken cancellationToken)
{
    // ... (cache-sjekk og semaphore)

    // Hent nytt token via OAuth2 password grant
    var response = await base.SendAsync(new HttpRequestMessage(HttpMethod.Post, $"{_config.BaseUrl}/Token")
    {
        Content = new FormUrlEncodedContent(new Dictionary<string, string>
            {
                ["grant_type"] = "password",
                ["username"] = _config.Username,
                ["password"] = _config.Password
            })
    }, cancellationToken).ConfigureAwait(false);
    response.EnsureSuccessStatusCode();

    // Deserialiser respons
    var result = await response.Content.ReadFromJsonAsync<NHNTokenResponse>
        (cancellationToken: cancellationToken).ConfigureAwait(false);

    cache.Token = result!.AccessToken;
    var serverExpiry = DateTime.UtcNow.AddSeconds(result.ExpiresIn - TokenExpiryBufferSeconds);
    cache.Expiry = serverExpiry < DateTime.UtcNow.AddDays(NHNTokenMaxAgeDays)
        ? serverExpiry
        : DateTime.UtcNow.AddDays(NHNTokenMaxAgeDays);

    return cache.Token;
}
\end{lstlisting}